%%------------------------------------------------------------------------------
%% General formatting
%%------------------------------------------------------------------------------

\newcommand{\keyword}[1]{{\texttt{#1}}}

\newcommand*{\numberinmargin}[1]{%
  \makebox[0pt][r]{%
    \makebox[\marginparwidth][l]{\arabic{#1}}%
    \hskip\marginparsep
  }%
}

\setcounter{secnumdepth}{3} % number subsections
% Only include sections (no subsections) in the TOC
\setcounter{tocdepth}{3}
% Ensure section entries are left-aligned (no extra indent)
\makeatletter
\renewcommand{\l@chapter}[2]{\@dottedtocline{0}{0pt}{2.6em}{\textbf{#1}}{#2}}
\renewcommand{\l@section}[2]{\@dottedtocline{1}{0pt}{2.6em}{\textbf{#1}}{#2}}
\makeatother

% Make spacing between paragraphs and titles consistent
\titleformat{\chapter}{\large\bfseries}{\thechapter.}{0em}{}
\titlespacing{\chapter}{0pt}{0.5\baselineskip}{0pt}


\titleformat{\section}{\bfseries}{\thesection.}{0em}{}
\titlespacing{\section}{0pt}{0.5\baselineskip}{0pt}

\titleformat{\subsection}{\bfseries}{\thesubsection.}{0em}{}
\titlespacing{\subsection}{0pt}{0.5\baselineskip}{0pt}

\titleformat{\subsubsection}{\bfseries}{\thesubsection.}{0em}{}
\titlespacing{\subsubsection}{0pt}{0.5\baselineskip}{0pt}

\newcommand{\nocontentsline}[3]{}
\let\origcontentsline\addcontentsline
\newcommand\stoptoc{\let\addcontentsline\nocontentsline}
\newcommand\resumetoc{\let\addcontentsline\origcontentsline}

\newcommand{\Ch}[2]{\vspace{1em}\chapter[#1]{#1\hfill[#2]}\label{#2}}
\newcommand{\ChUnnumbered}[2]{\vspace{1em}\chapter*{#1\hfill[#2]}\label{#2}\addcontentsline{toc}{chapter}{#1}}

\newcommand{\Sec}[2]{\section[#1]{#1\hfill[#2]}\label{#2}}
\newcommand{\Sub}[2]{\subsection[#1]{#1\hfill[#2]}\label{#2}}
\newcommand{\SubSub}[2]{\subsubsection[#1]{#1\hfill[#2]}\label{#2}}
\newcommand{\parnum}{\numberinmargin{parcount}}

% Don't start a new page for each chapter, and don't reset page numbers at the
% start of each chapter. Ths conforms to the Ecma house style (see: TOOLS-011,
% 6.2.1).
\titleclass{\chapter}{straight}

\newcommand{\newparcounter}[1]{
\newcounter{#1}
\counterwithin{#1}{chapter}
\counterwithin{#1}{section}
\counterwithin{#1}{subsection}
\counterwithin{#1}{subsubsection}
}

\newparcounter{parcount}
\newcommand\p{%
    \vspace{0.5\baselineskip}
    \stepcounter{parcount}%
    \parnum%
}

\newcommand{\Par}[2]{\paragraph[#1]{#1\hfill[#2]\\}\label{#2}\p}

% Table captions will be separated from the table number by a dash (see:
% TOOLS-011,6.2.5).
\captionsetup{labelsep=endash}

%%------------------------------------------------------------------------------
%% Grammar formatting
%%------------------------------------------------------------------------------

\newlength{\GrammarIndent}
\setlength{\GrammarIndent}{\leftmargini}
\newlength{\GrammarInc}
\setlength{\GrammarInc}{\GrammarIndent}
\newlength{\GrammarRest}
\setlength{\GrammarRest}{2\GrammarIndent}

\newenvironment{grammar} {
  \newcommand{\define}[1]{{\textit{##1}\textnormal{:}}}
  \newcommand{\terminal}[1]{{\texttt{##1}}}
  \newcommand{\br}{\\}
  \newcommand{\opt}[1]{{##1\textsubscript{\textit{opt}}}}

  \renewcommand{\texttt}[1]{{\small\ttfamily\upshape ##1}}

  \newcommand{\grammarindentfirst}{\GrammarIndent}
  \newcommand{\grammarindentinc}{\GrammarInc}
  \newcommand{\grammarindentrest}{\GrammarRest}
  \itshape

  \begin{grammarlist}
  \item\relax
}{
  \end{grammarlist}
}

\newlist{grammarlist}{itemize}{1}
\setlist[grammarlist]{
  parsep=1ex, partopsep=0pt, itemsep=0pt, topsep=0pt, label={},
  leftmargin=\grammarindentrest, listparindent=-\grammarindentinc,
  itemindent=\listparindent
}

\lstloadlanguages{C++} % TODO: Consider defining an HLSL language...
\lstnewenvironment{HLSL}
{\lstset{language=C++,
basicstyle=\small,
xleftmargin=1em}}{}

%%------------------------------------------------------------------------------
%% Note formatting
%%------------------------------------------------------------------------------

% Ecma style requires notes be numbered if more than one occurs in the same
% section. This doesn't currently happen in our spec, but we may need to define
% a note counter if it does in the future.
\newenvironment{note}
  {
   \vspace{0.5\baselineskip}
   \begin{tabular}{|p{3.0em} p{0.84\textwidth}|}
   \hline
   \begin{minipage}[t]{3.0em}\centering\textnormal{NOTE}
   \end{minipage} &
   \begin{minipage}[t]{0.84\textwidth}\raggedright
  }
  {
   \end{minipage} \\\hline
   \end{tabular}
   \vspace{0.5\baselineskip}
  }
